\exnumber{06}
\extitle{UX I - Ex 7 - Prototypage}
\exfiles{ex06/archive zip, nomme Name\_FirstName\_DeliverableName\_Date\_VersionNumber}
\exauthorize{Ceux nommés dans l'exercice}
\makeheaderfiles

\textbf{Instructions :}

Développez votre idée !

Concevez un storyboard de Pierre (ou de votre persona) utilisant votre idée comme solution. Montrez toutes les étapes (de 6 à 12) avec des dessins sur papier.

\textbf{Livrables :}

\begin{itemize}
\item Photos d'1 storyboard pour décrire votre solution, sur papier
\item Déposez les documents sur Github, dans un zip nommé « ProjectTitle\_Name\_FirstName », et chaque livrable ainsi : « Name\_FirstName\_DeliverableName\_Date\_VersionNumber ». Exemple : « Doe\_John\_MidFiPrototype\_05242024\_V1 ».
\item Ajoutez un titre dans tous les documents écrits.
\item Soignez l'aspect global des livrables. Ils doivent être clairs, simples et lisibles. Vous pouvez regarder canva.com, sans surcharger le dessin de détails.
\end{itemize}

\textbf{Pistes :}

\begin{itemize}
\item 1 dessin par action
\end{itemize}

\textbf{Ressources :}

\begin{itemize}
\item Comment les storyboards UX peuvent transformer votre processus créatif \url{https://uxdesign.cc/ux-storyboarding-2ce43875f3de?sk=41b28ecef8c52d9944e5e636881111cf}
\item 17 raisons d'utiliser les storyboards en UX design \url{https://uxdesign.cc/17-reasons-to-use-a-storyboards-in-ux-design-2bc6fea73e20?source=friends_link&sk=5bb2722f1f9ff06d4f22b0abfe52f508}
\item Comment storyboarder des expériences \url{https://uxdesign.cc/how-to-storyboard-experiences-fc051e2bc04d}
\item Utiliser des bandes dessinées et des storyboards pour tester vos concepts UX \url{https://uxdesign.cc/using-comic-strips-and-storyboards-to-test-your-ux-concepts-cccad7ac7f71}
\end{itemize}

\textbf{Citation :}

« Aucun produit n'est une île. Un produit est plus que le produit. C'est un ensemble cohérent et intégré d'expériences » Don Norman

